% ─── Chapter 9: Risk, Ethics & Considerations ─────────────────────
\chapter{Risk, Ethics \& Considerations}
\label{ch:risk}

\section{Risk Assessment}

\subsection{Technical Risks}

\begin{center}
\small
\begin{tabularx}{\textwidth}{l c X}
  \toprule
  \textbf{Risk} & \textbf{Impact} & \textbf{Mitigation} \\
  \midrule
  Medium-stage plateau &
    High &
    Systematic debugging through controlled experiments (positional
    encoding, embeddings, reward shaping).  Multiple architectural
    improvements have already been validated. \\
  Memory leaks &
    Medium &
    Observed RAM growth from 8\,GB to 52\,GB over 5~hours.  Mitigated with
    periodic cleanup, memory-leak counters, and the \texttt{memory\_safe}
    preset for limited hardware. \\
  RLlib API instability &
    Medium &
    The LSTM integration required workarounds for RLlib's connector path.
    Custom \texttt{RLModule} and \texttt{EnvRunner} implementations provide
    a stable interface. \\
  Reward hacking &
    Low &
    Multi-component reward with matchup and action-quality signals reduces
    the risk of the agent exploiting a single reward channel. \\
  \bottomrule
\end{tabularx}
\end{center}

\subsection{Timeline Risks}

The critical path depends on resolving the medium-stage plateau.  If the
reward signal redesign does not produce improvement within 2--3 experimental
iterations, the timeline for Milestone~1 completion could extend by an
additional 2--4~weeks.  This risk is accepted given the exploratory nature of
the project.

\section{Ethical Considerations}

\subsection{Domain Context}

This project operates entirely within the domain of \textbf{game AI}.  The
agent plays Pok\'{e}mon battles against simulated opponents in a controlled
environment.  There are no ethical concerns related to:

\begin{itemize}
  \item \textbf{Personal data:} No personal or sensitive data is collected,
        stored, or processed.  All data consists of game mechanics, trainer
        rosters, and battle states.
  \item \textbf{Harmful applications:} The trained model cannot be repurposed
        for harmful real-world applications.  It is specialised for a
        turn-based game with a fixed action space.
  \item \textbf{Bias or fairness:} The opponent populations are programmatically
        defined.  There are no human subjects, user-facing decisions, or
        societal impacts.
\end{itemize}

\subsection{Software Licences}

\begin{center}
\small
\begin{tabularx}{\textwidth}{l l}
  \toprule
  \textbf{Component} & \textbf{Licence} \\
  \midrule
  Pok\'{e}mon Showdown & AGPL-3.0 (open source, self-hosted) \\
  \texttt{poke-env}      & MIT \\
  Ray / RLlib           & Apache-2.0 \\
  PyTorch               & BSD-3-Clause \\
  MLflow                & Apache-2.0 \\
  \bottomrule
\end{tabularx}
\end{center}

All dependencies are open-source.  The project does not redistribute
copyrighted Pok\'{e}mon assets; it uses the Showdown simulator's built-in data
under fair use for research and educational purposes.

\subsection{EU AI Act Classification}

Under the EU AI Act, this system falls into the \textbf{minimal risk}
category.  It is a research prototype for a game environment with no
interaction with natural persons, no biometric data, no critical
infrastructure, and no social scoring.  No conformity assessment is required.

\section{Human-in-the-Loop Considerations}

Although the training loop is automated, several critical decisions require
human judgement:

\begin{itemize}
  \item \textbf{Checkpoint selection:} A human reviews MLflow dashboards to
        select the best checkpoint for gauntlet evaluation, rather than
        automatically selecting by training win rate (which can be misleading).
  \item \textbf{Curriculum stage overrides:} If a training run stagnates, a
        human may manually adjust the curriculum stage, opponent mix, or
        promotion thresholds.
  \item \textbf{Hyperparameter tuning:} Learning rate, entropy coefficient,
        and reward weights are set based on human interpretation of training
        curves---not automated tuning.
  \item \textbf{Experiment design:} The sequence of architectural changes
        (positional encoding, embeddings, action compression) was guided by
        human analysis of failure modes, not random search.
\end{itemize}

This human oversight is essential because the environment is complex enough
that automated optimisation could converge to local optima that do not
generalise to the gauntlet.

\section{Limitations}

\begin{itemize}
  \item The agent has been tested only on Milestone~1 (fixed team vs.\ random
        opponents).  Team builder integration (Milestone~3) and full gauntlet
        completion (Milestone~4) remain future work.
  \item Cross-turn memory via LSTM has been implemented but not yet validated
        end-to-end in a full training run.
  \item The system operates in the Gen\,8 random-battle format.  Transfer to
        other Pok\'{e}mon generations or formats would require observation-space
        adjustments.
  \item All training and evaluation use local hardware (no cloud deployment).
  \item The heuristic opponent (\texttt{SimpleHeuristicsPlayer}) is a
        rule-based baseline, not a strong competitive AI.  Success against it
        does not guarantee competence against human players or more
        sophisticated agents.
\end{itemize}
